\textbf{Phần 1. Trắc nghiệm lựa chọn}
\setcounterpageref{ex}{0}

\begin{ex} %[5D1N3]
Phân số nào dưới đây là phân số tối giản?
\choice
{\True $\dfrac{7}{9}$}
{$\dfrac{16}{18}$}
{$\dfrac{15}{12}$}
{$\dfrac{25}{30}$}
\loigiai{
Phân số $\dfrac{7}{9}$ có tử số và mẫu số không cùng chia hết cho số tự nhiên nào lớn hơn $1$, nên $\dfrac{7}{9}$ là phân số tối giản.\\
\textbf{Đáp án của câu hỏi là A}
}
\end{ex}

\begin{ex} %[5D1H3]
Rút gọn phân số $\dfrac{18}{24}$ về dạng tối giản ta được kết quả là
\choice
{$\dfrac{9}{12}$}
{$\dfrac{6}{8}$}
{\True $\dfrac{3}{4}$}
{$\dfrac{2}{3}$}
\loigiai{
Chia cả tử số và mẫu số cho $6$: $\dfrac{18}{24} = \dfrac{18 : 6}{24 : 6} = \dfrac{3}{4}$.\\
\textbf{Đáp án của câu hỏi là C}
}
\end{ex}

\begin{ex} %[5D1H3]
Rút gọn phân số $\dfrac{56}{42}$ về dạng phân số tối giản ta được
\choice
{$\dfrac{28}{21}$}
{\True $\dfrac{4}{3}$}
{$\dfrac{8}{6}$}
{$\dfrac{7}{6}$}
\loigiai{
Chia cả tử số và mẫu số cho $14$: $\dfrac{56}{42} = \dfrac{56 : 14}{42 : 14} = \dfrac{4}{3}$.\\
\textbf{Đáp án của câu hỏi là B}
}
\end{ex}

\begin{ex} %[5D1H3]
Số thích hợp điền vào dấu hỏi chấm để $\dfrac{24}{40} = \dfrac{12}{?} = \dfrac{3}{5}$ là
\choice
{\True $20$}
{$15$}
{$25$}
{$30$}
\loigiai{
Ta có $\dfrac{24}{40} = \dfrac{24 : 2}{40 : 2} = \dfrac{12}{20}$. Vậy số cần điền vào dấu hỏi chấm là $20$.\\
\textbf{Đáp án của câu hỏi là A}
}
\end{ex}

\begin{ex} %[5D1H3]
Quy đồng mẫu số hai phân số $\dfrac{9}{5}$ và $\dfrac{21}{40}$ với mẫu số chung nhỏ nhất là $40$, ta được hai phân số là
\choice
{$\dfrac{36}{40}$ và $\dfrac{21}{40}$}
{$\dfrac{45}{40}$ và $\dfrac{21}{40}$}
{$\dfrac{72}{40}$ và $\dfrac{42}{40}$}
{\True $\dfrac{72}{40}$ và $\dfrac{21}{40}$}
\loigiai{
Vì $40 : 5 = 8$ nên $\dfrac{9}{5} = \dfrac{9 \times 8}{5 \times 8} = \dfrac{72}{40}$. Giữ nguyên phân số $\dfrac{21}{40}$.\\
\textbf{Đáp án của câu hỏi là D}
}
\end{ex}

\begin{ex} %[5D1H3]
Tính giá trị của biểu thức $\dfrac{6 \times 12 \times 17}{17 \times 4 \times 20}$ bằng cách rút gọn ta được
\choice
{$\dfrac{3}{5}$}
{\True $\dfrac{9}{10}$}
{$\dfrac{18}{20}$}
{$1$}
\loigiai{
Ta có: $\dfrac{6 \times 12 \times 17}{17 \times 4 \times 20} = \dfrac{6 \times 3 \times 1}{1 \times 1 \times 20} = \dfrac{18}{20} = \dfrac{9}{10}$.\\
\textbf{Đáp án của câu hỏi là B}
}
\end{ex}

\begin{ex} %[5D1N3]
Phân số nào dưới đây lớn hơn $1$?
\choice
{$\dfrac{16}{17}$}
{$\dfrac{90}{100}$}
{$\dfrac{120}{125}$}
{\True $\dfrac{31}{29}$}
\loigiai{
Phân số có tử số lớn hơn mẫu số thì lớn hơn $1$. Vì $31 > 29$ nên $\dfrac{31}{29} > 1$.\\
\textbf{Đáp án của câu hỏi là D}
}
\end{ex}

\begin{ex} %[5D1H3]
Phân số nào dưới đây bé hơn phân số $\dfrac{4}{7}$?
\choice
{$\dfrac{5}{7}$}
{$\dfrac{20}{35}$}
{\True $\dfrac{12}{28}$}
{$\dfrac{13}{21}$}
\loigiai{
Ta có $\dfrac{4}{7} = \dfrac{16}{28}$. Phân số $\dfrac{12}{28} < \dfrac{16}{28} = \dfrac{4}{7}$.\\
(Lưu ý: $\dfrac{20}{35} = \dfrac{4}{7}$; $\dfrac{13}{21} > \dfrac{12}{21} = \dfrac{4}{7}$).\\
\textbf{Đáp án của câu hỏi là C}
}
\end{ex}

\begin{ex} %[5D1H3]
Phân số nào dưới đây lớn hơn phân số $\dfrac{5}{8}$?
\choice
{$\dfrac{10}{16}$}
{\True $\dfrac{25}{32}$}
{$\dfrac{13}{24}$}
{$\dfrac{3}{8}$}
\loigiai{
Quy đồng về mẫu số chung $32$: $\dfrac{5}{8} = \dfrac{20}{32}$. Ta thấy $\dfrac{25}{32} > \dfrac{20}{32}$.\\
\textbf{Đáp án của câu hỏi là B}
}
\end{ex}

\begin{ex} %[5D1H3]
Dấu thích hợp điền vào chỗ chấm để $\dfrac{7}{6} \dots \dfrac{29}{24}$ là
\choice
{$>$}
{$=$}
{\True $<$}
{Không so sánh được}
\loigiai{
Quy đồng mẫu số: $\dfrac{7}{6} = \dfrac{28}{24}$. Vì $\dfrac{28}{24} < \dfrac{29}{24}$ nên $\dfrac{7}{6} < \dfrac{29}{24}$.\\
\textbf{Đáp án của câu hỏi là C}
}
\end{ex}

\begin{ex} %[5D1V3]
Số học sinh khối 5 tham gia các môn thể thao ở Câu lạc bộ như sau: Cầu lông có $\dfrac{1}{3}$ số học sinh khối 5, Bóng đá có $\dfrac{2}{5}$, Võ có $\dfrac{1}{5}$, Bóng rổ có $\dfrac{4}{15}$. Môn thể thao được học sinh khối 5 tham gia nhiều nhất là
\choice
{Cầu lông}
{Võ}
{Bóng rổ}
{\True Bóng đá}
\loigiai{
Quy đồng mẫu số các phân số về mẫu số chung là $15$:\\
Cầu lông: $\dfrac{1}{3} = \dfrac{5}{15}$; Bóng đá: $\dfrac{2}{5} = \dfrac{6}{15}$; Võ: $\dfrac{1}{5} = \dfrac{3}{15}$; Bóng rổ: $\dfrac{4}{15}$.\\
Vì $\dfrac{6}{15} > \dfrac{5}{15} > \dfrac{4}{15} > \dfrac{3}{15}$ nên môn Bóng đá có số học sinh tham gia nhiều nhất.\\
\textbf{Đáp án của câu hỏi là D}
}
\end{ex}

\begin{ex} %[5D1V3]
Tìm số tự nhiên thích hợp điền vào dấu hỏi chấm để $\dfrac{3}{8} < \dfrac{?}{8} < \dfrac{4}{7}$.
\choice
{\True $4$}
{$2$}
{$3$}
{$5$}
\loigiai{
Quy đồng mẫu số chung là $56$:\\
$\dfrac{3}{8} = \dfrac{21}{56}$; $\dfrac{4}{7} = \dfrac{32}{56}$; $\dfrac{?}{8} = \dfrac{7 \times ?}{56}$.\\
Ta có: $21 < 7 \times ? < 32 \Rightarrow 7 \times ? = 28 \Rightarrow ? = 4$.\\
\textbf{Đáp án của câu hỏi là A}
}
\end{ex}

\textbf{Phần 2. Trắc nghiệm đúng sai}
\setcounterpageref{ex}{0}

\begin{ex} %[5D1H3]
Cho bốn phân số: $\dfrac{3}{4}$, $\dfrac{7}{12}$, $\dfrac{20}{16}$ và $\dfrac{19}{24}$. Xét tính đúng sai của các khẳng định sau:
\choiceTF
{\True Phân số $\dfrac{20}{16}$ sau khi rút gọn tối giản bằng $\dfrac{5}{4}$.}
{Phân số $\dfrac{7}{12}$ lớn hơn phân số $\dfrac{3}{4}$.}
{\True Trong bốn phân số đã cho, phân số lớn nhất là $\dfrac{20}{16}$.}
{\True Sắp xếp bốn phân số theo thứ tự từ bé đến lớn là: $\dfrac{7}{12}; \dfrac{3}{4}; \dfrac{19}{24}; \dfrac{20}{16}$.}
\loigiai{
Quy đồng mẫu số chung $48$ (hoặc $24$):\\
$\dfrac{7}{12} = \dfrac{14}{24}$; $\dfrac{3}{4} = \dfrac{18}{24}$; $\dfrac{19}{24}$; $\dfrac{20}{16} = \dfrac{5}{4} = \dfrac{30}{24}$.
\begin{itemize}
    \item Rút gọn: $\dfrac{20}{16} = \dfrac{20:4}{16:4} = \dfrac{5}{4}$.
    \item So sánh: $\dfrac{14}{24} < \dfrac{18}{24} \Rightarrow \dfrac{7}{12} < \dfrac{3}{4}$.
    \item Vì $\dfrac{20}{16} > 1$ còn ba phân số kia bé hơn $1$ nên $\dfrac{20}{16}$ là lớn nhất.
    \item Thứ tự từ bé đến lớn là: $\dfrac{14}{24} < \dfrac{18}{24} < \dfrac{19}{24} < \dfrac{30}{24}$, tức là $\dfrac{7}{12}; \dfrac{3}{4}; \dfrac{19}{24}; \dfrac{20}{16}$.
\end{itemize}
\begin{itemize}
    \item a) \textbf{Đúng}.
    \item b) \textbf{Sai}: $\dfrac{7}{12} < \dfrac{3}{4}$.
    \item c) \textbf{Đúng}.
    \item d) \textbf{Đúng}.
\end{itemize}
}
\end{ex}

\begin{ex} %[5D1V3]
Tại Câu lạc bộ thể thao trường Tiểu học Kim Đồng, số học sinh khối 5 tham gia: Cầu lông chiếm $\dfrac{1}{3}$, Bóng đá chiếm $\dfrac{2}{5}$, Võ chiếm $\dfrac{1}{5}$, Bóng rổ chiếm $\dfrac{4}{15}$ tổng số học sinh. Xét tính đúng sai của các khẳng định sau:
\choiceTF
{\True Số học sinh tham gia môn Võ ít nhất trong bốn môn.}
{Số học sinh tham gia môn Cầu lông nhiều hơn môn Bóng đá.}
{\True Tổng số phần học sinh tham gia môn Võ và Bóng đá là $\dfrac{3}{5}$ số học sinh khối 5.}
{Số học sinh tham gia môn Bóng rổ nhiều hơn môn Cầu lông.}
\loigiai{
Quy đồng về mẫu số $15$: Cầu lông là $\dfrac{5}{15}$; Bóng đá là $\dfrac{6}{15}$; Võ là $\dfrac{3}{15}$; Bóng rổ là $\dfrac{4}{15}$.
\begin{itemize}
    \item Môn Võ có tỉ lệ $\dfrac{3}{15}$ (nhỏ nhất) nên có ít học sinh tham gia nhất.
    \item Cầu lông ($\dfrac{5}{15}$) ít hơn Bóng đá ($\dfrac{6}{15}$).
    \item Tổng môn Võ và Bóng đá: $\dfrac{1}{5} + \dfrac{2}{5} = \dfrac{3}{5}$.
    \item Bóng rổ ($\dfrac{4}{15}$) ít hơn Cầu lông ($\dfrac{5}{15}$).
\end{itemize}
\begin{itemize}
    \item a) \textbf{Đúng}.
    \item b) \textbf{Sai}: Cầu lông ít hơn bóng đá.
    \item c) \textbf{Đúng}.
    \item d) \textbf{Sai}: Bóng rổ ít hơn cầu lông.
\end{itemize}
}
\end{ex}

\textbf{Phần 3. Trả lời ngắn}
\setcounterpageref{ex}{0}

\begin{ex} %[5D1H3]
Tìm số tự nhiên thích hợp điền vào dấu hỏi chấm để có đẳng thức:
\[
\dfrac{36}{54} = \dfrac{?}{9}.
\]
\shortans{6}
\loigiai{
\textbf{Đáp số:} $6$\\
Chia cả tử và mẫu của $\dfrac{36}{54}$ cho $6$: $\dfrac{36 : 6}{54 : 6} = \dfrac{6}{9} \Rightarrow ? = 6$.
}
\end{ex}

\begin{ex} %[5D1H3]
Rút gọn phân số $\dfrac{72}{90}$ thành phân số tối giản có dạng $\dfrac{a}{b}$. Tính giá trị của $a + b$.
\shortans{9}
\loigiai{
\textbf{Đáp số:} $9$\\
$\dfrac{72}{90} = \dfrac{72 : 18}{90 : 18} = \dfrac{4}{5}$. Do đó $a = 4, b = 5 \Rightarrow a + b = 9$.
}
\end{ex}

\begin{ex} %[5D1H3]
Tính giá trị của biểu thức:
\[
A = \dfrac{6 \times 12 \times 17}{17 \times 4 \times 20}.
\]
Viết kết quả dưới dạng phân số tối giản $\dfrac{m}{n}$. Tính tích $m \times n$.
\shortans{90}
\loigiai{
\textbf{Đáp số:} $90$\\
Rút gọn: $A = \dfrac{6 \times 3}{20} = \dfrac{18}{20} = \dfrac{9}{10}$. Khi đó $m = 9, n = 10 \Rightarrow m \times n = 90$.
}
\end{ex}

\begin{ex} %[5D1H3]
Tìm mẫu số chung nhỏ nhất của ba phân số $\dfrac{1}{3}$, $\dfrac{7}{12}$ và $\dfrac{35}{48}$.
\shortans{48}
\loigiai{
\textbf{Đáp số:} $48$\\
Vì $48$ chia hết cho cả $3$ và $12$ ($48 : 3 = 16$; $48 : 12 = 4$), nên mẫu số chung nhỏ nhất của ba phân số là $48$.
}
\end{ex}

\begin{ex} %[5D1V3]
Trong bốn phân số $\dfrac{3}{4}; \dfrac{7}{12}; \dfrac{20}{16}; \dfrac{19}{24}$, phân số nào có giá trị nhỏ nhất? (Điền tử số của phân số đó sau khi quy đồng mẫu số chung là $24$).
\shortans{14}
\loigiai{
\textbf{Đáp số:} $14$\\
Quy đồng mẫu số $24$: $\dfrac{7}{12} = \dfrac{14}{24}; \dfrac{3}{4} = \dfrac{18}{24}; \dfrac{19}{24}; \dfrac{20}{16} = \dfrac{30}{24}$. Phân số nhỏ nhất là $\dfrac{7}{12} = \dfrac{14}{24}$, tử số là $14$.
}
\end{ex}

\begin{ex} %[5D1V3]
Tìm chữ số tự nhiên $x$ thỏa mãn: $\dfrac{3}{8} < \dfrac{x}{8} < \dfrac{4}{7}$.
\shortans{4}
\loigiai{
\textbf{Đáp số:} $4$\\
Quy đồng mẫu số $56$: $\dfrac{21}{56} < \dfrac{7 \times x}{56} < \dfrac{32}{56}$. Suy ra $21 < 7 \times x < 32 \Rightarrow 7 \times x = 28 \Rightarrow x = 4$.
}
\end{ex}
